\begin{textsample}{NEG Dim 3 – 2010\_c4 – Score 1.00 – 2010\_c4/enade\_2010\_nutri\_97a31d8b0f.txt}  \label{ex:f3_neg_124}
CAPÍTULO 4 PERCEPÇÃO DA PROVA As análises feitas neste capítulo tratam das percepções dos estudantes, separados por grupos de Ingressantes e Concluintes, sobre a prova aplicada no ENADE/2010 na Área de Nutrição.
Estas percepções foram mensuradas por meio de nove questões que avaliaram desde o grau de dificuldade da prova até o tempo gasto para concluí -la.
As percepções sobre a prova foram relacionadas com o desempenho dos alunos e com a Grande Região de funcionamento do curso.
O desempenho dos alunos foi classificado em dois níveis : quarto inferior e quarto superior.
Para tanto, esse desempenho foi colocado em ordem ascendente.
O percentil 25, P25, também conhecido como primeiro quartil, é a nota de desempenho que deixa um quarto ( 25,0 \% ) dos valores observados abaixo e três quartos acima.
A Figura 1 apresenta uma ilustração destes conceitos.
O quarto inferior de desempenho é composto pelas notas abaixo do primeiro quartil.
Já o percentil 75, P75, também conhecido como terceiro quartil, é o valor para o qual há três quartos ( 75,0 \% ) dos dados abaixo e um quarto acima dele.
O quarto superior de desempenho é composto pelas notas iguais ou acima do terceiro quartil.
O percentil 50, P50, também conhecido como mediana, é o valor que divide as notas em dois conjuntos de igual tamanho.
Vale ressaltar que percentis, quartis e medianas são pontos, e não obrigatoriamente pertencentes ao conjunto original de dados.
Ao passo que os quartos são subconjuntos dos dados originais.
A seguir, serão apresentados Gráficos com resultados selecionados, relativos às nove questões avaliadas por grupos de estudantes.
As Tabelas no Anexo II apresentam os valores absolutos e a distribuição percentual das alternativas válidas das nove questões, segundo o mesmo recorte de desempenho dos alunos e Grande Região de funcionamento do curso. 4.1 GRAU DE DIFICULDADE DA PROVA 4.1.1 Componente de Formação Geral Na questão “ Qual o grau de dificuldade desta prova na parte de Formação Geral? ”, 20,5 \% do grupo de Ingressantes e 21,3 \% do grupo de Concluintes optaram pelas alternativas difícil ou muito difícil.
Entretanto, para mais da metade dos estudantes, tanto Ingressantes ( 64,8 \% ) quanto Concluintes ( 66,6 \% ), o Componente de Formação Geral da prova foi considerado com grau de dificuldade médio ( Gráfico 4.1 e Tabelas II.1 e II.10 no Anexo II ).
O percentual de Ingressantes que consideraram a prova difícil ou muito difícil foi mais elevado no quarto inferior ( 27,4 \% ) do que no superior ( 12,6 \% ), o que também foi verificado no grupo dos Concluintes, com percentual de participantes maior no quarto inferior ( 25,0 \% ) do que no quarto superior ( 15,4 \% ).
Semelhantemente ao ocorrido para os Ingressantes como um todo, para os quartos de menos e maior desempenho deste grupo, a alternativa modal para esta pergunta foi médio, com respectivamente 62,3 \% e 4.1.2 Componente de Conhecimento Específico Ao responderem à questão “ Qual o grau de dificuldade desta prova na parte de Componente Específico? ”, 53,2 \% do grupo de Ingressantes e 26,1 \% do grupo de Concluintes classificaram -na como difícil ou muito difícil.
O Componente de Conhecimento Específico da prova foi considerado com grau de dificuldade médio por 43,1 \% dos Ingressantes e por 66,0 \% dos Concluintes ( Gráfico 4.2 e Tabelas II.2 e II.11 no Anexo II ). 66,1 \% dos respondentes.
Já entre os Concluintes, estas percentagens foram, respectivamente, 65,7 \% e 67,9 \%.
Considerando -se as alternativas escolhidas pelos estudantes de cada Grande Região, a região Sul foi aquela onde houve maior proporção, no grupo de Ingressantes, com 25,4 \%, enquanto a de menor incidência foi a Nordeste, com 18,3 \% dos respondentes considerando a prova como difícil ou muito difícil.
Quanto ao grupo de Concluintes, que avaliou o Componente de Formação Geral da prova da mesma forma, 25,8 \% da região Sul tiveram, em maior medida, essa percepção da prova, enquanto 14,2 \% da região Nordeste tiveram a mesma percepção, em menor frequência.
Nas Grandes Regiões, a proporção de Ingressantes que consideraram o Componente de Formação Geral da prova como sendo de grau de dificuldade médio esteve localizada entre 62,7 \% na região Sul e 65,6 \% na região Nordeste.
A opção pela alternativa de grau médio por Grande Região que foi feita pelos Concluintes, atingiu a percentuais que estiveram entre 65,4 \% na Sul e 70,7 \%, na região Norte. 4.1.2 Componente de Conhecimento Específico Ao responderem à questão “ Qual o grau de dificuldade desta prova na parte de Componente Específico? ”, 53,2 \% do grupo de Ingressantes e 26,1 \% do grupo de Concluintes classificaram -na como difícil ou muito difícil.
O Componente de Conhecimento Específico da prova foi considerado com grau de dificuldade médio por 43,1 \% dos Ingressantes e por 66,0 \% dos Concluintes ( Gráfico 4.2 e Tabelas II.2 e II.11 no Anexo II ).
Considerando -se a percepção sobre o grau de dificuldade do Componente de Conhecimento Específico da prova de acordo com o desempenho dos estudantes, observa -se que 53,2 \% dos Ingressantes que obtiveram desempenho situado no quarto superior classificaram -na como difícil ou muito difícil.
Dos Ingressantes que obtiveram desempenho situado no quarto inferior, 51,2 \% optaram pela mesma alternativa.
Para os Concluintes, o percentual de estudantes que consideraram o Componente de Conhecimento Específico da prova difícil ou muito difícil foi um pouco maior entre os que obtiveram desempenho no quarto inferior ( 26,7 \% ) do que os que obtiveram desempenho no quarto superior ( 24,2 \% ).
Também entre os Concluintes, a alternativa modal para a Questão 2 foi o grau médio, com 65,2 \% do quarto inferior e 67,9 \% do superior.
Entre as Grandes Regiões brasileiras, a proporção de Ingressantes que avaliaram o Componente de Conhecimento Específico da prova como difícil ou muito difícil variou de 46,7 \% ( Norte ) a 56,2 \% ( Sudeste ).
O grupo de Concluintes avaliou esse componente como difícil ou muito difícil com menor frequência, com o mínimo alcançado na região Nordeste, 18,8 \%, e o máximo na região Sul, com 29,0 \% dos respondentes.
O percentual de Ingressantes que classificaram o grau de dificuldade como médio, no Componente de Conhecimento Específico da prova, variou de 40,4 \% ( regiões Sudeste e Sul ) a 50,3 \% ( Norte ).
Dentre os Concluintes, a percepção do mesmo grau de dificuldade aconteceu com maior frequência, variando de 63,4 \% ( Sul ) a 72,1 \% ( Nordeste ). 4.2 EXTENSÃO DA PROVA EM RELAÇÃO AO TEMPO TOTAL Indagados quanto à extensão da prova, em relação ao tempo total oferecido para a sua resolução, os estudantes apontaram, com maior incidência, a alternativa que considerava a extensão adequada, para ambos os grupos e todas as agregações consideradas ( Gráfico 4.3 e Tabelas II.3 e II.12 no Anexo II ).
No grupo de Ingressantes, a maioria, 61,0 \%, respondeu ser a extensão da prova adequada, ao passo que, no grupo de Concluintes, este valor foi menor : 49,4 \% tiveram a mesma opinião.
Para 35,3 \% dos Ingressantes o exame foi longo ou muito longo.
Essa mesma consideração foi feita por 47,5 \% dos Concluintes.
No grupo de Ingressantes, a maioria, tanto para o quarto de desempenho inferior ( 58,4 \% ) quanto para o de melhor desempenho ( 60,8 \% ) respondeu ser a extensão da prova adequada.
Já entre os Concluintes, esses valores foram menores, situados em 49,1 \% nos dois quartos.
Considerando -se ainda os quartos extremos de desempenho, observa -se que a proporção de estudantes que consideraram a prova longa ou muito longa, em relação ao tempo total destinado à sua resolução, acompanhou a percepção de Ingressantes e Concluintes.
Para o grupo de Ingressantes, entre os estudantes no quarto de desempenho inferior, 36,6 \% declararam ter esta percepção, ao passo que entre aqueles no quarto de melhor desempenho, esta proporção foi de 35,8 \%.
O mesmo pode ser constatado entre os Concluintes : 47,3 \% no quarto de menor desempenho e 47,9 \% no quarto de maior desempenho declararam ter tido esta percepção.
Na análise por Grandes Regiões, o percentual de Ingressantes que responderam ser a extensão da prova adequada variou de 59,3 \% no Nordeste a 65,9 \% no Norte.
Já entre os Concluintes, estes valores estavam compreendidos entre 49,0 \% ( Sudeste ) e 52,0 \% ( Centro-Oeste ).
Aqueles Ingressantes que avaliaram a prova como longa ou muito longa em relação ao tempo total destinado à sua resolução variou de 28,4 \% na região Norte a 37,8 \% na região Sul.
Os percentuais, quanto ao mesmo aspecto, entre os Concluintes, foram de 41,5 \% na região Centro-Oeste, indo até 48,7 \% na região Norte. 4.3 COMPREENSÃO DOS ENUNCIADOS DAS QUESTÕES 4.3.1 Componente de Formação Geral Com relação aos enunciados das questões do Componente de Formação Geral, as opiniões foram positivas, com percentuais que foram dos 76,5 \% de Concluintes aos 79,6 \% de Ingressantes que consideraram claros e objetivos os enunciados de todas ou da maioria das questões ( Gráfico 4.4 e Tabelas II.4 e II.13 no Anexo II ).
Ao se comparar as percepções dos estudantes segundo o desempenho, observou -se que, entre os Ingressantes do quarto superior, a clareza e objetividade de todos ou da maioria dos enunciados das questões foram percebidas por 86,1 \%.
A proporção apurada para os Concluintes desse mesmo quarto foi de 83,6 \%.
No quarto de desempenho inferior, tal avaliação foi emitida por 71,9 \% dos Ingressantes e por 67,2 \% dos Concluintes.
A partir dessas estatísticas, é possível afirmar que todos ou a maioria dos enunciados das questões do Componente de Formação Geral estavam claros e objetivos para a maioria dos respondentes, segundo os quartos de desempenho considerados, embora os que se situaram no quarto superior tenham declarado esta percepção com maior incidência.
Na análise regional, observa -se que as opiniões de que todos ou a maioria dos enunciados das questões do Componente de Formação Geral estavam claros e objetivos variaram, entre os Ingressantes, de 79,0 \% no Sudeste a 80,9 \% no Centro-Oeste.
Entre os Concluintes, os valores correspondentes variaram de 75,4 \% no Sudeste a 79,0 \% no Norte. 4.3.2 Componente de Conhecimento Específico Com relação aos enunciados das questões do Componente de Conhecimento Específico da prova, para os Concluintes da Área de Nutrição, a clareza e a objetividade foram percebidas em todos ou na maioria destes por 80,5 \%.
Para o grupo dos Ingressantes, foram percebidas por 74,9 \%.
A proporção de estudantes que considerou os enunciados das questões claros e objetivos foi mais elevada nos quartos superiores : 81,0 \% dos Ingressantes e 85,7 \% dos Concluintes, se comparados aos 68,2 \% e aos 73,2 \% dos mesmos grupos, respectivamente, que estiveram no quarto inferior de desempenho ( Gráfico 4.5 e Tabelas II.5 e II.14 no Anexo II ).
Nas Grandes Regiões brasileiras, o percentual de Concluintes que consideraram claros e objetivos todos ou a maioria dos enunciados das questões do Componente de Conhecimento Específico da prova variou entre 79,1 \% ( Norte ) e 82,7 \% ( Sul ).
Entre os Ingressantes, os índices variaram entre 72,1 \% na região Nordeste e 77,8 \% da Centro-Oeste, para estas mesmas respostas. 4.4 SUFICIÊNCIA DAS INFORMAÇÕES/INSTRUÇÕES FORNECIDAS Ao avaliarem as informações/instruções fornecidas para a resolução das questões, 81,2 \% do grupo dos Ingressantes e 89,5 \% do grupo dos Concluintes respondentes de todo o Brasil afirmaram que estas eram até excessivas ou suficientes em todas ou na maioria das questões ( Gráfico 4.6 e Tabelas II.6 e II.15 no Anexo II ).
Levando -se em conta o desempenho dos participantes, notam -se diferenças entre as opiniões de estudantes, Ingressantes e Concluintes, dos quartos inferior e superior de desempenho, como mostra o Gráfico 4.6.
O percentual de participantes que avaliaram as informações/instruções como excessivas ou suficientes em todas ou na maioria das questões foi mais elevado nos quartos superiores, correspondendo a 86,8 \% dos Ingressantes e a 94,5 \% dos Concluintes, percentuais superiores à média nacional para cada um dos dois grupos.
Já nos quartos inferiores, a suficiência das informações/instruções declarada como excessiva em todas ou na maioria das questões foi percebida por 75,1 \% dos Ingressantes e por 83,1 \% dos Concluintes.
Quanto à distribuição de respondentes nas Grandes Regiões, observa -se que as proporções de estudantes que perceberam como excessivas ou suficientes em todas ou na maioria das questões as informações/instruções fornecidas para a resolução das mesmas superaram as médias nacionais, em maior medida, entre os Ingressantes ( 84,0 \% ) e Concluintes ( 91,4 \% ) da região Sul. 4.5 DIFICULDADE ENCONTRADA AO RESPONDER À PROVA Perguntados sobre as dificuldades com as quais se depararam ao responder à prova, 63,2 \% do grupo de Ingressantes e 9,6 \% do grupo de Concluintes apontaram o desconhecimento do conteúdo.
Concluintes, por sua vez, apontaram a forma diferente de abordagem do conteúdo como dificuldade em maior evidência, mencionada por 45,8 \% deste grupo.
Esta alternativa foi apontada por uma parcela menor dos Ingressantes, 20,4 \%.
A falta de motivação para fazer a prova foi a dificuldade para 19,5 \% dos Concluintes e para 7,9 \% dos Ingressantes.
Considerando -se os respondentes em todo o Brasil, 7,0 \% dos Ingressantes e 19,6 \% dos Concluintes afirmaram que não tiveram qualquer tipo de dificuldade para responder à prova ( Tabelas II.7 e II.16 no Anexo II ).
Com relação aos quartos de desempenho, o desconhecimento do conteúdo foi a dificuldade modal entre os Ingressantes, apontada por 71,1 \% daqueles no quarto superior e por 53,3 \% dos do quarto inferior.
Entre os Concluintes, os percentuais de estudantes com a mesma percepção foi inferior ao encontrado entre os Ingressantes, 10,9 \% no quarto superior e 8,5 \% no quarto inferior.
A alternativa modal pelo grupo de Concluintes, quando desagregada pelo desempenho, foi que a dificuldade encontrada era a forma diferente de abordagem do conteúdo : 48,7 \% no quarto inferior e 40,3 \% do quarto superior assim o responderam.
Na análise por Grandes Regiões, o percentual de Ingressantes que apontaram o desconhecimento do conteúdo como dificuldade ao responder à prova variou de 57,6 \%, na região Norte a 65,1 \%, na Sudeste.
Entre os Concluintes, tal percentual foi mais baixo, variando de 7,4 \% na região Nordeste a 11,7 \%, na Sul.
A forma diferente de abordagem do conteúdo foi citada, pelos Concluintes, com percentuais que variaram entre 40,0 \% ( região Nordeste ) e 49,5 \% ( Centro-Oeste ).
Em contrapartida, os percentuais desta alternativa entre os Ingressantes foram mais baixos, variando de 19,9 \% ( região Sul ) a 22,8 \% ( região Norte ).
O percentual de Ingressantes que citou a falta de motivação como dificuldade superou a média nacional nas regiões Nordeste ( 8,8 \% ) e Centro-Oeste ( 9,2 \% ).
Entre os Concluintes, os percentuais para essa resposta variaram de 15,3 \% ( região Norte ) a 24,7 \% ( região Nordeste ).
Percentuais acima das médias nacionais dos Ingressantes das regiões Nordeste ( 8,4 \% ) e Norte ( 9,7 \% ) declararam que não tiveram qualquer dificuldade para responder à prova.
Entre os Concluintes, esta situação aconteceu nas regiões Sudeste ( 19,8 \% ) e Nordeste ( 20,6 \% ).
O Gráfico 4.7 apresenta os percentuais de estudantes, Ingressantes e Concluintes, que apontaram o desconhecimento do conteúdo como dificuldade percebida ao responder à prova. 4.6 CONTEÚDOS DAS QUESTÕES OBJETIVAS DA PROVA Ao analisarem os conteúdos das questões objetivas da prova, a maior parte dos Ingressantes, 73,6 \% do total nacional, afirmou que não estudou ainda a maioria desses conteúdos, enquanto entre os Concluintes, apenas 1,3 \% expressaram a mesma opinião.
Por outro lado, a maioria dos Concluintes ( 81,9 \% ) afirmou ter estudado e aprendido muitos ou todos os conteúdos.
Para os Ingressantes, o percentual correspondente a essas alternativas foi 13,8 \% ( Gráfico 4.8 e Tabelas II.8 e II.17 no Anexo II ).
Considerando -se separadamente as opiniões de estudantes dos quartos de desempenho inferior e superior, observa -se que, para os Ingressantes, o percentual de participantes que afirmaram que não estudaram ainda a maioria desses conteúdos, foi mais elevado no quarto superior ( 74,3 \% ) do que no quarto inferior ( 70,3 \% ).
Para os Concluintes do quarto inferior, 3,3 \% ofereceram essa resposta, sendo menos de 1,0 \% os do quarto superior.
Tendo em conta os quartos superiores para os dois grupos, 89,6 \% dos Concluintes e 16,5 \% dos Ingressantes afirmaram ter estudado e aprendido muitos ou todos os conteúdos.
Na análise por Grande Região, a proporção de Ingressantes que afirmou que não estudou ainda a maioria desses conteúdos superou a média nacional nas regiões Nordeste ( 75,1 \% ) e Sul ( 75,4 \% ).
Já entre os Concluintes, as regiões Sudeste ( 1,4 \% ) e Norte ( 3,7 \% ) apresentaram percentuais maiores do que a média nacional para esta alternativa.
Em todas as regiões, a maioria dos Concluintes afirmaram ter estudado e aprendido muitos ou todos os conteúdos, variando entre 74,7 \% na região Norte e 83,4 \% na Nordeste. 4.7 TEMPO GASTO PARA CONCLUIR A PROVA Ao responderem sobre o tempo de conclusão da prova, mais da metade dos estudantes afirmou ter gasto entre duas e quatro horas.
Considerando -se separadamente as percepções do grupo de Ingressantes e do grupo de Concluintes, os percentuais, em termos nacionais, foram de 61,4 \% e 75,1 \%, respectivamente ( Gráfico 4.9 e Tabelas II.9 e II.18 no Anexo II ).
Uma vez sendo analisadas as alternativas oferecidas pelos estudantes que se situaram nos quartos superior e inferior, observa -se que uma maior proporção de Ingressantes no quarto superior declarou ter gasto entre duas e quatro horas para concluir a prova quando comparados com os do quarto inferior, respectivamente 76,6 \% e 47,0 \%.
A mesma situação foi verificada no grupo dos Concluintes, no qual 82,6 \% daqueles no quarto superior e 64,3 \% daqueles no quarto inferior declaram ter gasto entre duas e quatro horas para concluir a prova.
Nas cinco Grandes Regiões brasileiras, os Ingressantes que utilizaram entre duas e quatro horas para finalizar a prova superaram o percentual nacional nas regiões Nordeste ( 62,4 \% ), Sudeste ( 62,7 \% ) e Centro-Oeste ( 66,2 \% ).
Entre os Concluintes, o mesmo foi verificado entre os estudantes da região Sudeste ( 77,2 \% ).
As regiões Norte e Sul foram as que apresentaram o menor percentual de Concluintes que dispensaram entre duas e quatro horas para concluir a prova, ambas com 72,3 \%, como mostra o Gráfico 4.9.

% matched lemmas: 
\end{textsample}
